% Tipo de documento. En este caso es un art�culo, para folios A4, tama�o de la fuente 11pt y con p�gina separada para el t�tulo
\documentclass[a4paper,12pt,titlepage]{article}

% Carga de paquetes necesarios. OrdenesArticle es un paquete personalizado
\usepackage[spanish]{babel} 
\RequirePackage[T1]{fontenc}
\RequirePackage[ansinew]{inputenx} 
\usepackage[spanish,cap,cont,title,fancy]{OrdenesArticle}
\usepackage{array}
\usepackage{graphicx}
\usepackage{hyperref}
\usepackage{pifont}
\usepackage{listings}
\usepackage[usenames,dvipsnames]{color}
\usepackage{colortbl}
\usepackage{makeidx}
\hypersetup{bookmarksopen,bookmarksopenlevel=3,linktocpage,colorlinks,urlcolor=blue,citecolor=blue,
						linkcolor=blue,filecolor=blue,pdfnewwindow,
						pdftitle={Listado de preguntas Funcionalidad 1}, 
						pdfauthor={Juan Andrada Romero, Juan Jos� Antequera Flores},
						pdfsubject={Modelos Avanzados de Base de Datos}}


% Macro para definir una lista personalizada 
\newenvironment{milista}%
{\begin{list}{\textbullet}%
{\settowidth{\labelwidth}{\textbullet} \setlength{\leftmargin}{\dimexpr\labelsep+\labelwidth+5pt}
\setlength{\itemsep}{\dimexpr 0.5ex plus 0.25ex minus 0.25ex}
\setlength{\parsep}{\itemsep}
\setlength{\partopsep}{\itemsep}
\addtolength{\topsep}{-7.5pt}
}}%
{\end{list}}


\begin{document}

% En las p�ginas de portada e �ndices, no hay encabezado ni pie de p�gina
\pagestyle{empty} 

% Se incluye la portada
\thispagestyle{empty}
\begin{center}
  {\LARGE UNIVERSIDAD DE CASTILLA-LA MANCHA} \\
  \bigskip
  {\Large ESCUELA SUPERIOR DE INFORM�TICA} \\
  \vspace{28mm}
  \includegraphics[scale=0.45, keepaspectratio]{./imagenes/esi_bw.png} \\
  \vspace{30mm}
  {\LARGE \textbf{Modelos Avanzados de Base de Datos}} \\
  \vspace{14mm}
  {\Large \textsf{\textsc{Listado de preguntas del trabajo Distribuci�n 1:}}} \\
  \vspace{2mm}
  {\Large \textsf{\textsc{Bases de Datos Distribuidas}}}\\
  \vspace{15mm}
  {\large Juan Andrada Romero} \\
  {\large Juan Jos� Antequera Flores} \\
  {\large Jose Domingo L�pez L�pez} \\
  \vspace{9mm}
  {\large \today}
\end{center}
\clearpage

\parskip = 10pt
% Comienza el contenido del documento. Se utilizan n�meros ar�bigos y el encabezado y pie de p�gina personalizado
\pagenumbering{arabic}
\pagestyle{fancy}

\section{Preguntas seleccionadas}
\paragraph{Pregunta:} En el trabajo se comentan los lenguajes de definici�n de objetos y de consulta de objetos que define el est�ndar ODMG, pero no se hace referencia a ning�n lenguaje para poder manipular objetos, es decir, insertarlos en la base de datos, modificarlos o eliminarlos. �Dicho est�ndar define un lenguaje de manipulaci�n de datos espec�fico? �C�mo se realizar�an las manipulaciones sobre objetos en una base de datos orientada a objetos?

\paragraph{Respuesta:} El est�ndar ODMG-93 sugiere que el lenguaje de manipulaci�n de objetos (OML) sea la extensi�n de un lenguaje de programaci�n, de forma que se pueden realizar entre otras las siguientes operaciones sobre la base de datos: creaci�n, borrado, y modificaci�n de un objeto. \cite{bdoo}.

Por tanto, los objetos se manipulan dentro del lenguaje de programaci�n orientado a objetos, utilizando el concepto de transacci�n y de clases persistentes. De este modo, como indica el ODMG API (\cite{odmgapi}), para realizar las operaciones de inserci�n, modificaci�n o eliminaci�n de un objeto dentro del lenguaje de programaci�n, se proceder�a de la siguiente forma:
\begin{enumerate}
	\item Crear una transacci�n.
	\item Obtener la base de datos en uso.
	\item Ejecutar la operaci�n de inserci�n, modificaci�n o eliminaci�n.
\end{enumerate}

Un ejemplo de inserci�n ser�a el siguiente:

\begin{verbatim}
 public static void storeNewProduct(Product product)
	    Transaction tx = odmg.newTransaction();
	    tx.begin();
	    // get current used Database instance
	    Database db = odmg.getDatabase(null);
	    // make persistent new object
	    db.makePersistent(product);
	    tx.commit();
\end{verbatim}

\newpage
\paragraph{Pregunta:} En el documento se comentan algunas caracter�sticas de los SGBDOR, pero no se mencionan las de un SGBDOO. �Cu�les son las caracter�sticas que debe un cumplir un SGBDOO para que sea considerado como tal?

\paragraph{Respuesta:} En 1989 se hizo el Manifiesto \cite{manifest} de los sistemas de base de datos orientados a objetos el cual propuso trece caracter�sticas obligatorias para un SGBDOO y cuatro opcionales, las cuales son:
\begin{milista}
\item Caracter�sticas obligatorias de orientaci�n a objetos:
\begin{enumerate}
\item Deben soportarse objetos complejos
\item Deben soportarse mecanismos de identidad de los objetos
\item Debe soportarse la encapsulaci�n
\item Deben soportarse los tipos o clases
\item Los tipos o clases deben ser capaces de heredar de sus ancestros
\item Debe soportarse el enlace din�mico
\item El DML debe ser computacionalmente complejo
\item El conjunto de todos los tipos de datos debe ser ampliable
\end{enumerate}
\item Caracter�sticas obligatorias de SGBD:
\begin{enumerate}
\item Debe proporcionarse persistencia a los datos
\item El SGBD debe ser capaz de gestionar bases de datos de muy gran tama�o
\item El SGBD debe soportar a usuarios concurrentes
\item El SGBD debe ser capaz de recuperarse de fallos hardware y software
\item El SGBD debe proporcionar una forma simple de consultar los datos.
\end{enumerate}
\item Caracter�sticas opcionales:
\begin{enumerate}
\item Herencia m�ltiple
\item Comprobaci�n de tipos e inferencia de tipos
\item Sistema de base de datos distribuido
\item Soporte de versiones
\end{enumerate}
\end{milista}\newpage
\section{Preguntas descartadas}
\paragraph{Pregunta:} Si las bases de datos orientadas a objetos surgieron para combatir las limitaciones de las bases de datos relacionales, �por qu� siguen siendo las bases de datos relacionales las m�s utilizadas?

\paragraph{Respuesta:} Se siguen usando porque es un est�ndar muy arraigado y extendido, y en muchas empresas es muy costoso hacer un cambio a nuevas tecnolog�as aunque aporten m�s ventajas que las antiguas. Adem�s, el est�ndar de las bases de datos orientadas a objetos no est� todo lo maduro que debiera para poder hacer la migraci�n sin correr ning�n riesgo.
Por otra parte, las bases de datos orientadas a objetos suelen utilizarse en aplicaciones de dise�o y fabricaci�n por computador (CAD), que no est�n tan extendidas como otras aplicaciones comerciales que hacen uso de bases de datos relacionales. Por esta raz�n, se siguen utilizando en mayor medida, por estar m�s extendidas en las aplicaciones comerciales. 
\newpage
\paragraph{Pregunta:} �Qu� ventajas ofrece un sistema gestor de bases de datos orientadas a objetos frente a un sistema cl�sico relacional? \cite{ventajas}

\paragraph{Respuesta:} 
\begin{milista}
\item Los objetos en la OODMBS pueden almacenar un n�mero arbitrario de tipos at�micos, as� como otros objetos.
\item Los datos en el mundo real tienen habitualmente caracter�sticas jer�rquicas.
\item No es necesario un lenguaje de consultas para acceder a datos de la OODB, al contrario que ocurre con RDB, ya que la interacci�n con la base de datos es transparente al acceder a los objetos.
\item En una aplicaci�n t�pica que usa un lenguaje de programaci�n orientado a objetos y un RDMS se necesita bastante tiempo en "`mapear"' las tablas a partir de los objetos.
\item No hay claves primarias.
\item Un modelo de datos deber�a modelar entidades y sus relaciones, restricciones y operaciones que cambian el estado de los datos en el sistema.
\end{milista}

\newpage
% Bibliograf\'ia 
\bibliographystyle{plain}
\pagestyle{plain} 
\bibliography{Bibliografia}


\end{document}

